\section{Main Theorem}

\begin{theorem}[Conditional quotient-first private PAC bound]\label{thm:main}
Assume Assumptions~\ref{assump:finite-littlestone},
\ref{assump:countable-evaluation-quotient},
\ref{assump:realizable-iid}, and
\ref{assump:approximate-dp-regime}.  Let $K_C^{\mathrm{VC\text{-}Lyu}}$
be the totalized quotient-first law defined in the preliminaries, and let
\[
 A_N^{\mathrm{VC}}(s,E):=K_C^{\mathrm{VC\text{-}Lyu}}(T_N(s),E).
\]
If $d=0$, then $N=0$, the output is the unique member of $\bar C$, and
$A_0^{\mathrm{VC}}$ is $(0,0)$-DP with zero decoded population error.  If
$d\ge1$, there are universal constants $K_{\mathrm V}\ge1$ and
$q_{\mathrm V}\in\mathbb N_0$ such that $A_N^{\mathrm{VC}}$ is a measurable
Markov kernel, is $(\varepsilon,\delta)$-DP for every raw replace-one pair
(including nonrealizable labels), and satisfies
\[
 \sup_D\sup_{c\in C}
 \Pr_{S\sim P_{D,c}^{N},\,\bar H\sim K_C^{\mathrm{VC\text{-}Lyu}}(T_N(S),\cdot)}
 \left[\operatorname{err}_D(\operatorname{Dec}_C(\bar H),c)>\alpha\right]
 \le\beta,
\]
with
\[
 N\le K_{\mathrm V}\Lambda^{q_{\mathrm V}}R_{\mathrm{VC}}.
\]
Here $(d,v,\alpha,\beta,\varepsilon,\delta)$ are the exposed variables;
$K_{\mathrm V}$ and $q_{\mathrm V}$ are universal and may depend on none of
$(X,\Sigma,C,D,c)$ or on any decomposition, list, partition, event, or
output.  The probability mode is unconditional iid high probability over
sampling, partition, and mechanism randomness; the horizon mode is fixed
sample size $N$; and the norm is the population binary zero-one error.

On the same quotient interface, the totalized old-Lyu law has a measurable
raw pullback, is all-input $(\varepsilon,\delta)$-DP, and has the same
realizable PAC guarantee with
\[
 N_{\mathrm{old}}\le K_{\mathrm O}\Lambda^{q_{\mathrm O}}R_{\mathrm{old}}
\]
for universal $K_{\mathrm O},q_{\mathrm O}$.  When $|C|<\infty$, the finite
class law has the same output, decoder, privacy, and PAC interfaces with
\[
 N_{\mathrm{fin}}\le K_{\mathrm F}\Lambda^{q_{\mathrm F}}R_{\mathrm{fin}};
\]
when $|C|=\infty$, this arm is assigned cost $+\infty$ and no finite
surrogate is substituted.  Consequently, choosing the least certified arm
gives a single learner (not a mixture of laws) satisfying
\[
 m_C(\alpha,\beta;\varepsilon,\delta)
 \le K_*\Lambda^{q_*}\min\{R_{\mathrm{fin}},R_{\mathrm{old}},R_{\mathrm{VC}}\},
\]
where all constants and exponents are universal and the minimum uses the
$+\infty$ convention above.
\end{theorem}

\begin{corollary}[Rate specialization and baseline frontier]\label{cor:frontier}
Under the theorem assumptions, specialize the auxiliary teacher and block
parameters by the least feasible choices proved in the appendix, retain the
ceiling terms, and use the explicit source-compatible schedule
$\delta K\Lambda^{q_*}R_{\mathrm{VC}}\to0$.  The resulting fixed-parameter
comparison is
\[
 m_C=\widetilde O\!\left(\min\{\log^+|C|,d^5,vd^4\}\right).
\]
The auxiliary choices, verification of every technical condition, displayed
term-absorption inequalities, probability conversion, and hidden-constant
dependence are proved in Proposition~\ref{prop:step-016-frontier}.
The exact baseline reductions are: $N=0$ when $d=0$; $vd^4=d^5$ when
$v=d$ (and the same polynomial scale when $v=\Theta(d)$); the finite arm is
available only for finite $C$; and no conclusion for uncountable evaluation
quotients or a universal polynomial in $v$ and $\log d$ is asserted.
\end{corollary}

